\subsection{Clocking and power management}

\begin{frame}
  \frametitle{Enabling hardware}
  In order to become functional, hardware blocks require:
  \begin{itemize}
  \item Power to be applied
  \item Clocks to be ticking
    \begin{itemize}
    \item Hardware blocks are connected through hardware buses (AHB, APB,
      AXI...), their interface logic require an input clock to respond
      \begin{itemize}
      \item One will experience bus hangs otherwise
      \end{itemize}
    \item When a device exposes a bus (eg. a SPI controller), there is
      typically a second clock whose frequency can be tuned
    \end{itemize}
  \item Resets to be released
  \end{itemize}
\end{frame}

\begin{frame}
  \frametitle{Power handling}
  \begin{itemize}
  \item Internal and external devices can be fed by a regulator
    \begin{itemize}
    \item Shared refcounted resources
    \item See \kfunc{devm_regulator_get_enable} 
    \end{itemize}
  \item Internal devices can typically be part of a power management domain
    (named \code{pmdomain}, formerly \code{genpd})
    \begin{itemize}
    \item Shared refcounted resources
    \item For example: all display related controllers may be in one power
      domain (CRTC, LCD controller, HDMI phy, etc). The whole domain
      is either powered on or off.
    \end{itemize}
  \end{itemize}
\end{frame}

\begin{frame}
  \frametitle{Clocks handling}
  \begin{itemize}
  \item A clock tree typically starts from the main crystal, feeds PLLs, gates,
    divisors and reaches every device in the system
    \begin{itemize}
    \item Shared refcounted resources
    \item Managed by the Common Clock Framework (CCF)
    \item Simple API described in \kfile{include/linux/clk.h}.
      \begin{itemize}
      \item \kfunc{devm_clk_get} to lookup and obtain a reference to a clock producer
      \item \kfunc{clk_prepare_enable} to inform the system when the clock source should be running
      \item \kfunc{clk_disable_unprepare} to inform the system when the clock source is no longer required.
      \item \kfunc{clk_get_rate} to obtain the current clock rate (in Hz) for a clock source
      \item \kfunc{clk_set_rate} to set the current clock rate (in Hz) of a clock source
      \end{itemize}
    \end{itemize}
  \item Allows to declare the available clocks and their association
    to devices in the Device Tree
  \item Provides a {\em debugfs} representation of the clock tree
  \item Is implemented in \kdir{drivers/clk}
  \end{itemize}
\end{frame}

\begin{frame}{Diagram overview of the common clock framework}
  \begin{center}
    \includegraphics[width=\textwidth]{slides/kernel-power-management/clock-framework.pdf}
  \end{center}
\end{frame}

\begin{frame}
  \frametitle{Reset handling}
  \begin{itemize}
  \item Goal is to put a device in a known harder default state
  \item Typically done while probing
  \item Resets can be simple GPIOs
  \item More complex resets can be registered in the reset control framework
    \begin{itemize}
    \item Pointed by the \code{resets} DT property
    \item API is straightforward:
      \begin{itemize}
      \item \kfunc{devm_reset_control_get}
      \item \kfunc{reset_control_assert}
      \item \kfunc{reset_control_deassert}
      \end{itemize}
    \end{itemize}
  \end{itemize}
\end{frame}

\begin{frame}
  \frametitle{Runtime Power Management}
  \begin{itemize}
  \item The state of a device can change at runtime to save power
    \begin{itemize}
    \item Sort of per-device idle state
    \end{itemize}
  \item According to the kernel configuration interface: \emph{Enable
      functionality allowing I/O devices to be put into energy-saving
      (low power) states at run time (or autosuspended) after a
      specified period of inactivity and woken up in response to a
      hardware-generated wake-up event or a driver's request.}
  \item New hooks must be added to the drivers:
    \code{runtime_suspend()}, \code{runtime_resume()},
    \code{runtime_idle()} in the \kstruct{dev_pm_ops}
    structure of \kstruct{device_driver}.
  \item API and details on \kdochtml{power/runtime_pm}:
    \begin{itemize}
    \item \kfunc{pm_runtime_enable}
    \item \kfunc{pm_runtime_get_sync}
    \item \kfunc{pm_runtime_put}
    \end{itemize}
  \item Runtime PM may use resets, clocks, power...
    \begin{itemize}
    \item It sometimes replaces direct clock lookups and handling on some old
      ARM platforms, like the AM335x.
    \end{itemize}
  \end{itemize}
\end{frame}

\begin{frame}
  \frametitle{Useful resources}
  \begin{itemize}
  \item \kdochtmldir{power} in kernel documentation.
  \item Introduction to kernel power management \\
    Kevin Hilman (Kernel Recipes 2015)
    \begin{itemize}
    \item \url{https://www.youtube.com/watch?v=juJJZORgVwI}
    \end{itemize}
  \item Overview of Generic PM Domains \\
    Kevin Hilman (Kernel Recipes 2017)
    \begin{itemize}
    \item \url{https://www.youtube.com/watch?v=SctfvoskABM}
    \end{itemize}
  \item Linux Power Management Features, Their Relationships and
    Interactions \\
    Théo Lebrun (Embedded Linux Conference Europe 2024)
    \begin{itemize}
    \item \url{https://www.youtube.com/watch?v=_jb6U40ZCZk}
    \end{itemize}
  \end{itemize}
\end{frame}
